%!TEX root = ../../../main.tex

\chapter{Test- und Evaluationsstrategie}
\label{chap:evaluation}

\paragraph{Testsätze (Leitfrage~6).} Die KPI-Suite (\texttt{src/eval}) bewertet vier Sätze mit denselben Regeln. Der \emph{Kernsatz} enthält die neun Pflichtfragen: fünf Standardfragen, zwei Grenzfälle (Vorauszahlung Wasser gegen Vertrag, Verträge mit Verlängerung um ein Jahr) und zwei Fragen ohne Beleg. Der \emph{Entwicklungssatz} (54 Fragen) ist der einzige, an dem ich Fehler analysiert und Einstellungen gewählt habe. Der \emph{synthetische Satz} (51 Fragen) entsteht per Vorlage aus den Sollwerten, ähnlich der Testset-Generierung von RAGAS \cite{es2024}. Der \emph{zurückgehaltene Satz} (20 Fragen, alle fünf Typen) war vor dem Umbau geschrieben und per SHA-256 versiegelt.

\paragraph{Bewertung.} Statt eines LLM-Richters wie in RAGAS \cite{es2024}, der lokal langsam und selbst unzuverlässig wäre, prüfe ich deterministisch gegen Sollwerte. Korrekt ist eine Antwort, wenn jeder Sollwert in einem nicht verneinten Satz steht. Als \emph{Halluzination} zählt eine unbelegte Zahl oder eine Antwort auf eine unbeantwortbare Frage. Die Maße folgen der Literatur: Zitation nach ALCE \cite{gao2023alce}, Recall@5 und \gls{MRR} nach \cite{manning2008}, Enthaltung nach \cite{wen2024}, Werkzeugwahl nach \cite{qin2024}, Effizienz und Robustheit nach HELM \cite{liang2023}. Jede der 23 KPIs hat in \texttt{kpis.py} eine Definition, einen Zielwert und eine Begründung. Die Intervalle sind Wilson- \cite{wilson1927} bzw.\ Bootstrap-Intervalle \cite{efron1993}. Jede Antwort wird mit Trajektorie als JSONL gespeichert, personenbezogene Daten maskiert. Version1 lief unverändert auf denselben Fragen (Tabelle~\ref{tab:kpis}).

% automatisch erzeugt von src/eval/latex.py
\begin{table}[tp]
\centering
\scriptsize
\setlength{\tabcolsep}{3pt}
\begin{tabularx}{\textwidth}{@{}l l Y Y Y Y@{}}
\toprule
\textbf{KPI} & \textbf{Ziel} & \textbf{V1 P1} & \textbf{V1 P2} & \textbf{V2 P1} & \textbf{V2 P2} \\
\midrule
\multicolumn{6}{@{}l}{\textit{Qualität}} \\
Korrektheit & $\geq$ 80 \% & 53\,\% [37--68] & 53\,\% [37--68] & 48\,\% [36--61] & 52\,\% [39--64] \\
Treue & $\geq$ 95 \% & 96\,\% [91--99] & 98\,\% [95--100] & 98\,\% [96--100] & 82\,\% [69--92] \\
Halluzinationsrate & $\leq$ 5 \% & 16\,\% [7--30] & 11\,\% [4--24] & 10\,\% [5--21] & 22\,\% [14--35] \\
Zitationsrate & $\geq$ 90 \% & 90\,\% [70--97] & 89\,\% [67--97] & 94\,\% [80--98] & 52\,\% [37--67] \\
Zitationspräzision & $\geq$ 85 \% & 65\,\% [43--82] & 67\,\% [44--84] & 88\,\% [73--95] & 52\,\% [37--67] \\
Enthaltung & $\geq$ 90 \% & 86\,\% [49--97] & 86\,\% [49--97] & 60\,\% [31--83] & 60\,\% [31--83] \\
Fehl-Enthaltung & $\leq$ 10 \% & 35\,\% [21--53] & 42\,\% [26--59] & 31\,\% [20--45] & 17\,\% [9--30] \\
Recall@5 & $\geq$ 90 \% & 40\,\% [26--56] & 40\,\% [26--56] & 85\,\% [76--93] & 85\,\% [76--93] \\
Precision@5 & $\geq$ 60 \% & 43\,\% [26--59] & 43\,\% [26--59] & 72\,\% [61--82] & 72\,\% [61--82] \\
MRR & $\geq$ 0,80 & 0,49 [0,33--0,65] & 0,49 [0,33--0,65] & 0,86 [0,77--0,95] & 0,86 [0,77--0,95] \\
Beleg@5 & $\geq$ 85 \% & 52\,\% [34--69] & 52\,\% [34--69] & 77\,\% [63--87] & 77\,\% [63--87] \\
Feldgenauigkeit der Extraktion & $\geq$ 95 \% & 49\,\% [41--58] & 49\,\% [41--58] & 96\,\% [91--98] & 96\,\% [91--98] \\
\multicolumn{6}{@{}l}{\textit{Agent}} \\
Tool-Selection-Accuracy & $\geq$ 85 \% & -- & 10\,\% [3--25] & -- & 75\,\% [61--85] \\
Schritte je Frage & $\leq$ 4 & 0,0 [0,0--0,0] & 0,7 [0,4--1,0] & 0,0 [0,0--0,0] & 1,3 [1,1--1,4] \\
\multicolumn{6}{@{}l}{\textit{Effizienz}} \\
Latenz p50 & $\leq$ 60 s & 105\,s [92--136] & 176\,s [64--292] & 28\,s [11--38] & 36\,s [30--57] \\
Latenz p95 & $\leq$ 180 s & 254\,s [173--1876] & 948\,s [550--2935] & 62\,s [50--76] & 134\,s [105--195] \\
Tokens/s & $\geq$ 5 & 4,2 [3,9--4,4] & 3,5 [3,3--3,7] & 5,5 [5,4--5,5] & 5,1 [5,0--5,2] \\
Spitzen-RAM & $\leq$ 12 GB & 8,2\,GB & 8,2\,GB & 6,6\,GB & 9,1\,GB \\
Ingest-Zeit je Dokument & $\leq$ 120 s & -- & -- & 154\,s [140--171] & 154\,s [140--171] \\
Rewriting-Latenz & $<$ 500 ms & 0\,ms [0--0] & 0\,ms [0--0] & 2776\,ms [2546--2919] & 2776\,ms [2546--2919] \\
\multicolumn{6}{@{}l}{\textit{Robustheit}} \\
Konsistenz bei Wiederholung & $\geq$ 90 \% & 100\,\% [70--100] & 100\,\% [70--100] & 100\,\% [88--100] & 100\,\% [88--100] \\
Fehlerquote & $\leq$ 2 \% & 0\,\% [0--9] & 0\,\% [0--9] & 0\,\% [0--6] & 0\,\% [0--6] \\
Ausfallquote & 0 \% & 0\,\% [0--9] & 0\,\% [0--9] & 0\,\% [0--6] & 0\,\% [0--6] \\
\bottomrule
\end{tabularx}
\caption{KPIs der Ausgangsmessung (Version1) und der Schlussmessung (Version2) auf denselben Fragen, Punktwert mit 95-\%-Intervall (Anteile: Wilson, Mittelwerte und Perzentile: Bootstrap)}
\label{tab:kpis}
\end{table}


\paragraph{Ergebnis.} Version2 beantwortet nicht messbar mehr Fragen richtig als Version1: 48\,\% (P1) und 52\,\% (P2) gegen 53\,\%. Messbar besser ist alles vor dem Sprachmodell: Recall@5 stieg von 40 auf 85\,\%, die MRR von 0,49 auf 0,86, die Feldgenauigkeit von 49 auf 96\,\%, die richtige Werkzeugwahl von 10 auf 75\,\%. Die Latenz im Median sank von 105 bzw.\ 176\,s auf 28 bzw.\ 36\,s. P1 halluziniert seltener (10 statt 16\,\%) und zitiert präziser (88 statt 65\,\%). Im Kernsatz löst P2 16 von 18 Antworten, im zurückgehaltenen Satz lösen beide Prototypen nur 14 von 40. Auf dem Entwicklungssatz erreicht P1 nach der Abstimmung 81\,\%, auf ungesehenen Fragen 35\,\%. Die Abstimmung hat also vor allem Formulierungen der bekannten Fragen gelernt.

\paragraph{P1 gegen P2.} P1 löst Tabellen- und Rechenfragen besser (12 gegen 4 von 18), weil es die Zeilen direkt im Kontext sieht. P2 liest sie über \texttt{search\_tables} falsch, etwa 8300 statt 5500\,kWh. Fristfragen löst P2 vollständig (12 von 12), weil \texttt{list\_contracts} sie berechnet. P1 löst davon nur 2 und lehnt bei 31\,\% der beantwortbaren Fragen ab. P2 halluziniert dafür häufiger (22\,\%): Nach einem Gas- oder Fernwärmevertrag gefragt, antwortet es mit dem Heizstromvertrag.

\paragraph{Weiche zwischen P1 und P2.} Eine feste Textregel schickt Fragen zu Fristen, Listen über mehrere Dokumente und Vergleiche an P2, alle übrigen an P1 (\texttt{router.py}, ohne Sprachmodell). Die Idee stammt aus der Fehleranalyse oben, also aus dem zurückgehaltenen Satz. Die Muster und die Annahmeregel R9 wurden nur auf dem Entwicklungssatz festgelegt. Dort brachte die Weiche nichts: 44 von 54 Antworten richtig, genau wie P1 allein. R9 verlangt nur Gleichstand, daher wurde sie übernommen und einmal auf Kern- und zurückgehaltenem Satz ausgewertet, aus den Antworten der Schlussmessung. Ergebnis: 59\,\% richtig (34 von 58) gegen 48 und 52\,\%. Gegenüber P1 sind das $+10$ Prozentpunkte mit Intervall $[0;\,+21]$, gegenüber P2 $+7$ mit $[-7;\,+22]$, also nicht gesichert. Die Halluzinationsrate liegt mit 17\,\% zwischen beiden, die Latenz im Median bei 38\,s. Weil der zurückgehaltene Satz damit zum dritten Mal gesehen wurde und die Idee aus ihm stammt, ist der Wert eher optimistisch.

\paragraph{Transparenz.} Die erste Schlussmessung zeigte zwei Fehler. Die Bewertung erkannte die Ablehnungsformel von Version2 nicht; korrigiert, wurden alle Läufe neu bewertet. Außerdem gab P1 in 30 von 58 Antworten nur eine Belegnummer aus. Ich habe den Prompt auf dem Entwicklungssatz korrigiert und P1 neu gemessen. Für P1 wurde der zurückgehaltene Satz damit zweimal verwendet. Beide Läufe liegen im Repository.
